\wysubtitle{logos.pdf}{出题：\textsf{LJL \& CTZ \& ZWX}}
{https://qm.qq.com/q/OlkSiy9W04}{HDU 数学营}
{http://weixin.qq.com/r/mhKYgFTEMQlOrRDj90eI}{未央学社公众号}
\begin{flushleft}
    \kaishu\qquad 本资料仅作为模拟练习之用，目的是为了帮助大家更有效地复习，并减轻对考试的担忧。请正确的对待此资料，其旨在辅助复习，而非预示具体的考试内容。我们鼓励同学们认真复习，大学学习主打理解，而非刷题，期望大家在期中考试中取得优异成绩。
\end{flushleft}

\section{选择题}
\begin{choice}{B}
    以下命题正确的是
\begin{tasks}(2)
    \task 无界数列必趋于无穷大
    \task 趋于无穷大的数列必无界
    \task 趋于正无穷大的数列必在充分大时单调
    \task 不趋于无穷大的数列必有界
\end{tasks}
\end{choice}
\begin{solution}
\begin{itemize}
    \item 设$f(x)=\frac{1}{x^2}\sin\frac1x$，$\forall M>0,\ \exists x_0=\frac{1}{2[M]\pi}\in \overset{\circ}{U}(0),\ \text{s.t.}f(x_0)=0$.
    \item $\exists x_0'=\frac{1}{2[M]\pi+\pi/2}\in\overset{\circ}{U}(0)$，s.t. $f(x_0')=(2[M]\pi+\pi/2)^2>M$.
    \item 由无穷大量和无界量的定义可知当$x\to 0$时，$f(x)$为无界量，但不是无穷大量.
\end{itemize}
\end{solution}

\begin{choice}{B}
    $f(x)=\lim_{t\to 0}\ab(1+\frac{\sin t}{x})^{x^2/t}$在$(-\infty,+\infty)$内
\begin{tasks}(4)
    \task 连续
    \task 有可去间断点
    \task 有跳跃间断点
    \task 有无穷间断点
\end{tasks}
\end{choice}
\solute1
\begin{solution}*
    求函数相关性质时首先考虑定义域. $f(x)=\lim_{t\to 0}\ab(1+\frac{\sin x}{x})^{\frac{x}{\sin t}\cdot\frac{\sin t}{t}\cdot x}=\e^x$定义域为$(-\infty,0)\cup (0,+\infty)$.
\end{solution}

\begin{choice}{D}
    函数$f(x)=\arctan\frac{1}{1-x}$，当$x\to 1$时的极限值为
\begin{tasks}(4)
    \task $\frac\pi2$
    \task $-\frac\pi2$
    \task $0$
    \task 不存在
\end{tasks}
\end{choice}
\solute{.5}
\begin{solution}*
    求极限值时，首先判断极限值是否存在. $\lim_{x\to 1^+}f(x)=-\frac\pi2\neq\lim_{x\to 1^-}f(x)=\frac\pi2$，所以不存在极限.
\end{solution}

\begin{choice}{C}
    $F(x)=\e^x-(ax^2+bx+c)$至多有几个零点
\begin{tasks}(4)
    \task 1
    \task 2
    \task 3
    \task 4
\end{tasks}
\end{choice}
\solute2
\begin{solution}
\begin{itemize}
    \item 假设$F(x)$有4个零点（$x_0<x_1<x_2<x_3$），由罗尔定理得：在$(x_0,x_1)$存在$\xi_1\in(x_0,x_1)$，s.t. $F'(\xi_1)=0$.
    \item 同理，在$(x_1,x_2)$和$(x_2,x_3)$分别存在$\xi_2\in(x_1,x_2),\ \xi_3\in(x_2,x_3)$，s.t. $F'(\xi_2)=F'(\xi_3)=0$.
    \item 由于$F(\xi_1)=F'(\xi_2)=F'(\xi_3)=0$，所以在$(\xi_1,\xi_2)$中存在$\xi_4\in (\xi_1,\xi_2)$，s.t. $F''(\xi_4)=\e^{\xi_4}-2a=0$.
    \item 同理，存在$\xi_5\in(\xi_2,\xi_3)$，s.t. $F''(\xi_5)=\e^{\xi_5}-2a=0$.
    \item 由于$F''(\xi_4)=F''(\xi_5)=0$，所以存在$\xi_6\in(\xi_4,\xi_5)$，s.t. $F'''(\xi_6)=0$. 而$F'''(x)=\e^x\neq 0$，所以假设不成立.
    \item 综上所述，$F(x)$至多有3个零点.
\end{itemize}
\end{solution}

\begin{choice}{D}
    若$f(x+1)=af(x)$总成立，且$f'(0)=b$，其中$a,b$为非零常数，则$f(x)$在$x=1$
\begin{tasks}(4)
    \task 不可导
    \task 可导，且$f'(1)=a$
    \task 可导，且$f'(1)=b$
    \task 可导，且$f'(1)=ab$
\end{tasks}
\end{choice}
\solute1
\begin{solution}*这道题考查主要在于对函数导数定义的理解，可从定义出发去求解. 由题意
\begin{itemize}
    \item $f'(0)=\lim_{\Delta x\to 0}\frac{f(0+\Delta x)-f(0)}{\Delta x}=\lim_{\Delta x\to 0}\frac{f(\Delta x)-f(0)}{\Delta x}=b$
    \item $f'(1)=\lim_{\Delta x\to 0}\frac{f(1+\Delta x)-f(1)}{\Delta x}\xlongequal[f(1+\Delta x)=af(\Delta x)]{f(1+0)=af(0)}\lim_{\Delta x\to 0}\frac{af(\Delta x)-af(0)}{\Delta x}=ab$.
\end{itemize}    
\end{solution}

\begin{choice}{A}
    设函数$f(x)=\begin{cases}
        x^\alpha\cos\frac{1}{x^\beta} & x>0\\
        0 & x\leq 0
    \end{cases}\ (\alpha>0,\ \beta>0)$，若$f'(x)$在$x=0$处可导，则
\begin{tasks}(4)
    \task $\alpha-\beta>1$
    \task $0<\alpha-\beta\leq 1$
    \task $\alpha-\beta>2$
    \task $0<\alpha-\beta\leq 2$
\end{tasks}
\end{choice}
\solute1
\begin{solution}*
由于$f'(x)$在$x=0$处可导
\begin{itemize}
    \item $f'(0)=\lim_{x\to 0}\frac{f(x)-f(0)}{x}=\lim_{x\to 0}x^{\alpha -1}\cos\frac{1}{x^\beta}$存在，所以$\alpha-1>0,\ \alpha>1$.
    \item $f'(x)$在$x=0$处连续，因为$f'(0)=0$，所以
    \[\lim_{x\to 0}f'(x)=f'(0),\ \lim_{x\to 0}\alpha\cdot x^{\alpha-1}\cos\frac{1}{x^\beta}+\beta x^{\alpha-\beta-1}\sin\frac{1}{x^\beta}\Longrightarrow\alpha-\beta-1>0,\ \alpha-\beta>1\]
\end{itemize}
\end{solution}

\begin{choice}{C}
    设连续函数$f(x)$满足$\int_0^xf(t)\d t=f(x)-1$，则$f(x)=$
\begin{tasks}(4)
    \task $y=2\e^x$
    \task $y=\e^{2x}$
    \task $y=\e^x$
    \task $y=\e^{x^2}$
\end{tasks}
\end{choice}
\solute1
\begin{solution}*
    令$x=0$时，$\int_0^0f(t)\d t=f(0)-1=0\Rightarrow f(0)=1$. 对原式两边进行求导得
    \[f'(x)=f(x)\Rightarrow y'=y\Rightarrow\d y/y=\d x\]
    
    所以该微分方程通解为$y=C\e^x$. 将$x=0$代入得$y=1$，故结果为$f(x)=\e^x$.
\end{solution}

\begin{choice}{C}
    设常数$k>0$，函数$f(x)=\ln x-x/\e+k$在$(0,+\infty)$内零点个数为
\begin{tasks}(4)
    \task 0
    \task 1
    \task 2
    \task 3
\end{tasks}
\end{choice}
\solute1
\begin{solution}
\begin{itemize}
    \item 显然$f(x)$为初等函数，在$(0,+\infty)$连续. $\lim_{x\to0^+}f(x)=f(0^+)<0,\ \lim_{x\to +\infty}f(x)<0,\ f(\e)=k>0$.
    \item 由$f'(x)=1/x-1/\e$得$f(x)$在$(0,\e)$单调递增，$(\e,+\infty)$单调递减.
    \item 由于$f(0^+)\cdot f(\e)<0,\ f(\e)\cdot\lim_{x\to+\infty}f(x)<0$，所以$f(x)$在$(0,+\infty)$内有两个零点.
\end{itemize}
\end{solution}

\section{填空题}

\begin{problem}
    设函数$f(x)=(\e^x-b)/[(x-a)(x-1)]$有无穷间断点$x=0$及可去间断点$x=1$. 则$a=$\ans{$0$}，$b=$\ans{$\e$}.
\end{problem}
\begin{solution}
\begin{itemize}
    \item 已知$\lim_{x\to 0}f(x)=\infty$，则$\lim_{x\to 0}1/f(x)=0,\ a/(1-b)=0,\ a=0,\ b\neq 1$.
    \item 因为$f(x)$存在$x=1$的可去间断点，所以$\lim_{x\to 1}f(x)$存在.
    \item 由于$\lim_{x\to 1}x(x-1)=0$，所以$\lim_{x\to 1}(\e^x-b)=0,\ b=\e$.
\end{itemize}
\end{solution}

\vskip-2ex
\begin{paracol}{2}
\begin{problem}
求定积分$\int_0^{2\pi}(t-\sin t)^2\sin t\d t=$\ans{$-6\pi^2$}.
\end{problem}
\switchcolumn
\begin{solution}*
    定积分遇到幂函数与三角函教结合可用换元让上下限关于原点对称，从而可利用偶倍奇零.
\end{solution}
\end{paracol}
\vspace{-.75em}
\solute{1.5}
\begin{solution}*
    令$t=u+\pi$，原式
    \[\begin{aligned}
        I&=\int_0^{2\pi}(t^2\sin t-2t\sin^2t+\sin^3t)\d t=\int_{-\pi}^\pi[-(u^2+2\pi u+\pi^2)\sin u-2(u+\pi)\sin^2u-\sin^3u]\d u\\
        &=-4\pi\underset{\text{分部积分}}{\uwave{\int_0^\pi u\sin u\d u}}-4\pi\underset{\text{点火公式}}{\uwave{\int_0^\pi\sin^2u\d u}}=-6\pi^2
    \end{aligned}\]
\end{solution}

\begin{problem}
    求微分方程$\d y/\d x=x(y-1)$的通解为\ans{$x^2/2-\ln(y-1)+C=0$}.
\end{problem}
\solute0
\begin{solution}*
    原式可化为$\frac{\d y}{y-1}=x\d x\Rightarrow\int\frac{\d y}{y-1}=\int x\d x,\ \frac{x^2}{2}-\ln(y-1)+C=0$.
\end{solution}

\vskip-2ex
\begin{paracol}{2}
\begin{problem}
    求极限$\lim_{x\to 1}\frac{1-x^2}{\sin\pi x}$.
\end{problem}
\switchcolumn
\begin{solution}*
    求极限常用的手段为等价无小，所以我们自然想到要把极限趋于零，因此采用换元手段.
\end{solution}
\end{paracol}
\vspace{-.75em}
\solute0
\begin{solution}*
    令$t=x-1$，原式$I=\lim_{t\to 0}\frac{-(2+t)t}{\sin(1+t)\pi}=\lim_{t\to 0}\frac{t(t+2)}{\sin\pi t}=\lim_{t\to 0}\frac{2+t}{\pi}=\frac2\pi$.
\end{solution}

\section{计算题}
\begin{problem}
    求积分$\int\sec x\d x$.
\end{problem}
\solute2
\begin{solution}
    \[\begin{aligned}
        \int\sec x\d x&=\int\frac{1}{\cos x}\d x=\int\frac{\cos x}{\cos^2 x}\d x=\int\frac{\d\sin x}{1-\sin^2x}=\frac12\int\ab(\frac{1}{1-\sin x}+\frac{1}{1+\sin x})\d\sin x\\
        &=\frac12(-\ln\abs{1-\sin x}+\ln\abs{1+\sin x})+C=\frac12\ln\abs{\frac{1+\sin x}{1-\sin x}}+C=\frac12\ln\abs{\frac{(1+\sin x)(1+\sin x)}{(1-\sin x)(1+\sin x)}}+C\\
        &=\frac12\ln\abs{\frac{(1+\sin x)^2}{\cos^2x}}+C=\ln\abs{\frac{1+\sin x}{\cos x}}+C=\ln\abs{\frac{1}{\cos x}+\frac{\sin x}{\cos x}}+C=\ln\abs{\sec x+\tan x}+C
    \end{aligned}\]
\end{solution}

\begin{problem}
    求极限$\lim_{x\to 0}\frac{\e-\ln\ab[x+\exp(1+x)^{1/x}]}{x}$.
\end{problem}
\solute2
\begin{solution}*本题反复使用拉格朗日中值定理求极限.
\[\begin{aligned}[t]
    I&=\lim_{x\to 0}\frac{\ln\e^\e-\ln\ab[x+\exp(1+x)^{1/x}]}{x}=\lim_{x\to 0}\frac{\e^\e-[x+\exp(1+x)^{1/x}]}{x}\cdot\frac{1}{\e^\e}=\lim_{x\to 0}\frac{\e^\e-\exp(1+x)^{1/x}}{x}\cdot\frac{1}{\e^\e}-\frac{1}{\e^\e}\\
    &=\lim_{x\to 0}\frac{\e-(1+x)^{1/x}}{x}\cdot\e^\e\cdot\frac{1}{\e^\e}-\frac{1}{\e^\e}=\lim_{x\to 0}\frac{\e-(1+x)^{1/x}}{x}-\frac{1}{\e^\e}=\lim_{x\to 0}\frac{\e-\exp\ab[\dfrac{\ln(1+x)}{x}]}{x}-\frac{1}{\e^\e}\\
    &=\lim_{x\to 0}\frac{\ab[1-\dfrac{\ln(1+x)}{x}]}{x}\cdot\e-\frac{1}{\e^\e}=\lim_{x\to 0}\frac{[x-\ln(1+x)]}{x^2}\cdot\e-\frac{1}{\e^2}=\lim_{x\to 0}\frac{x^2/2+O(x^2)}{x^2}\cdot\e-\frac{1}{\e^\e}=\frac\e2-\frac{1}{\e^\e}
\end{aligned}\]
\end{solution}

\begin{problem}
    $y''+6y'+10y=0$，求通解
\end{problem}
\solute1
\begin{solution}*
    特征方程为$r^2+6r+10=0$，$r=-3\pm\i$. $y=\e^{-3x}(C_1\cos x+C_2\sin x)$.
\end{solution}

\begin{problem}
    $x\e^x-y\e^y=\sin(xy)$. 求$y'$和$y'\Big|_{x=0}$.
\end{problem}
\solute1
\begin{solution}
\begin{itemize}
    \item 等式两侧同时求导$-y'(\e^y+y\e^y)-y'x\cos(xy)=y\cos(xy)-\e^x(1+x)$.
    \item 得$y'=\frac{\e^x(1+x)-y\cos(xy)}{\e^x+y\e^y+x\cos(xy)}$. 当$x=0,\ y=1$时,$y'\Big|_{x=0}=\frac{1-0}{1}=1$.
\end{itemize}
\end{solution}

\begin{problem}
    $\int\frac{11\sin x-\cos x}{\sin x+\cos x}\d x$
\end{problem}
\solute1
\begin{solution}
    \[I=\int\frac{5(\sin x+\cos x)-6(\cos x+\sin x)}{\sin x+\cos x}\d x=\int 5\d x-6\int\frac{\d(\sin x+\cos x)}{\sin x+\cos x}=5x-6\ln\abs{\sin x+\cos x}+C\]
\end{solution}

\begin{problem}
    求定积分$\int_0^x(x-t)^2\e^{-2t}\d t$.
\end{problem}
\solute2
\begin{solution}
\begin{itemize}
    \item $F(x)=\int_0^xt^2\e^{-2x+2t}\d t=\int_0^x\frac{t^2\cdot\e^{2t}}{\e^{2x}}\d t=\e^{-2x}\int_0^x t^2\e^{2t}\d t$.
    \item $\int t^2\e^{2t}\d t=\frac12t^2\e^{2t}-\frac12\int t\d\e^{2t}=\frac12t^2\e^{2t}-\ab(\frac12t\e^{2t}-\frac12\int\e^{2t}\d t)=\frac12t^2\e^{2t}-\frac12t\e^{2t}+\frac14\e^{2t}$.
    \item $F(x)=\e^{-2x}\ab(\frac12x^2\e^{2x}-\frac12x\e^{2x}+\frac14\e^{2x}-\frac14)=\frac{x^2}{2}-\frac x2+\frac14-\frac14\e^{-2x}$.
\end{itemize}
\end{solution}

\vskip-2ex
\begin{paracol}{2}
\begin{problem}
    已知平面区域$D$由抛物线$y=x^2/5$、其在$\ab(5/2,5/4)$处的切线、直线$x=5$以及$x$轴围成.
\begin{enumerate}
    \item 平面区域的面积.
    \item $D$绕$y=x$旋转一周旋转体体积.
\end{enumerate}
\end{problem}
\switchcolumn\centering
\vfill
\begin{tikzpicture}[scale=.64]
    \filldraw [pattern=north west lines,2号色!80,fill=2号色!40,very thick,variable=\x,domain=0:5,samples=200,smooth,line cap=round] plot ({.6*\x},{.6*\x^2/5}) --++ (0,-.75) --++ (-2.25,-2.25) -- (0,0) -- cycle;
    \draw [thick,->] (-1,0) -- (4,0) node [anchor=south] {$x$};
    \draw [thick,->] (0,-1) -- (0,4) node [anchor=east] {$y$};
    \draw [very thick,line cap=round] (-.5,-.5) -- (3.5,3.5) node [anchor=south] {$y=x$};
    \draw[very thick,1号色,variable=\x,domain=0:5.5,samples=200,smooth,line cap=round] plot ({.6*\x},{.6*\x^2/5});
    \draw [dashed,yshift=-.75cm,thick,line cap=round] (.25,.25) -- (3.5,3.5);
    \draw [dashed,thick,line cap=round] (1.5,3.5) -- (1.5,-.5);
    \node [anchor=north] at (.75,-.5) {$\scriptstyle 5/4$};
    \node [anchor=north] at (1.5,-.5) {$\scriptstyle 5/2$};
    \node [anchor=north] at (3,-.5) {$\scriptstyle 5$};
    \draw [dashed,thick,line cap=round] (3,3.5) -- (3,-.5);
    \node [anchor=west] at (3,3) {$\scriptstyle\ab(5,\tfrac{15}{4})$};
\end{tikzpicture}
\vfill
\end{paracol}
\vspace{-.75em}
\solute3
\begin{solution}
\begin{enumerate}
    \item $y'=\frac25x,\ y'\Big|_{x=5/2}=1$，所以在$\ab(\frac52,\frac54)$处切线为$y-\frac54=x-\frac52,\ y=x-\frac54$.
    
    $y=\frac{x^2}{5}$与$y=x$的交点为$(0,2),\ (5,5)$；$y=x-\frac54$与$y=0,\ x=5$的交点为$\ab(\frac54,0),\ \ab(5,\frac{15}{4})$.
    \[S_D=\int_0^5\frac{x^2}{5}\d x-\int_{5/4}^5\ab(x-\frac54)\d x=\frac{125}{96}\]
    \item $V=\pi\ab(\frac{5/4}{\sqrt2})^2\cdot\frac{5/4}{\sqrt2}\cdot\frac13\cdot2+\pi\ab(\frac{5/4}{\sqrt2})^2\cdot\frac{15}{4}\sqrt2-\sqrt2\int_0^5\frac{\pi(x^2/5-x^2)}{2}\d x=\frac{225\sqrt2\pi}{192}$.
\end{enumerate}
\end{solution}

\begin{problem}
    求积分$\int_0^4\frac{\d x}{\sqrt{x(4-x)}}$.
\end{problem}
\solute1
\begin{solution}
    \[\int_0^4\frac{\d x}{\sqrt{x(4-x)}}=\int_0^4\frac{1}{\sqrt{4-(x-2)^2}}\d x=\int_0^4\frac{1}{\sqrt{1-\ab(\frac{x-2}{2})^2}}\d\ab(\frac{x-2}{2})=\arcsin\ab(\frac{x-2}{2})\Bigg|_0^4=\pi\]
\end{solution}

\begin{problem}
    计算$y=\frac{\sqrt[2]{x^2+x}\sqrt[3]{x+2}}{\sqrt[3]{x^2+2x}\sqrt[5]{x^3+2}}$的导函数.
\end{problem}
\solute2
\begin{solution}
\begin{itemize}
    \item $\ln y=\frac12\ln(x^2+x)+\frac13\ln(x+2)-\frac13\ln(x^2+2x)-\frac15\ln(x^3+2)$.
    \item $\frac{y'}{y}=\frac12\cdot\frac{2x+1}{x^2+x}+\frac13\cdot\frac{1}{x+2}-\frac13\cdot\frac{2x+2}{x^2+2x}-\frac15\cdot\frac{3x^2}{x^3+2}$.
    \item $y'=\ab(\frac12\cdot\frac{2x-1}{x^2+x}+\frac13\cdot\frac{1}{x+2}-\frac13\cdot\frac{2x+2}{x^2+2x}-\frac15\cdot\frac{3x^2}{x^3+2})y$.
\end{itemize}
\end{solution}

\begin{problem}
    微分方程：$f''(x)+af'(x)+f(x)=0\ (a>2)$. $f(0)=m,\ f'(0)=n$，求$\int_0^{+\infty}f(x)\d x$.
\end{problem}
\solute1
\begin{solution}
\begin{itemize}
    \item 特征方程$r^2+ar+1=0$有实根$r_{1,2}=\frac{-a\pm\sqrt{a^2-4}}{2}<0$.
    \item 所以$f(x)=C_1\e^{r_1x}+C_2\e^{r_2x},\ f'(x)=C_1r_1\e^{r_1x}+C_2r_2\e^{r_2x},\ \lim_{x\to\infty}f(x)=\lim_{x\to\infty}f'(x)=0$.
    \item 原式$I=\int_0^{+\infty}[-f''(x)-af'(x)]\d x=[-f'(x)-af(x)]\Big|_0^{+\infty}=f'(0)+af(0)=n+ma$.
\end{itemize}
\end{solution}

\begin{problem}
    设$f(x)$满足$f(1)=1,\ x\in\ab[1,\sqrt{\e^2-1}],\ f'(x)=\frac{2x}{x^2+f^2(x)}$，证明$f(x)\leq 3-\ln 2$.
\end{problem}
\solute1
\begin{solution}
\begin{proof}*
\begin{itemize}
    \item $f'(x)=\frac{2x}{x^2+f^2(x)}>0$，所以$f(x)$在$[1,\sqrt{\e^2-1}]$单调递增，$f'(x)\leq\frac{2x}{x^2+f^2(1)}=\frac{2x}{x^2+1}$.
    \item 所以$\int_1^x f'(x)\d x=f(x)-1\leq\int_1^x\frac{2x}{x^2+1}\d x,\ f(x)\leq 3-\ln 2$.\QED
\end{itemize}
\end{proof}
\end{solution}


\vfill
\rightline{/*Powered by \pkg{LiteSolution} \textsf\LaTeX~Template. Designed by WeYoung \href{https://ctan.org/home/xiamyphys}{Axia}.*/}